\documentclass[10pt,a4paper]{article}
\usepackage[margin=22mm]{geometry}
\usepackage{amsmath,booktabs}
\usepackage[hidelinks]{hyperref}
\setlength{\parindent}{0pt}
\setlength{\parskip}{5pt}
\setlength{\emergencystretch}{2em}
\title{\vspace{-1.8em}\textbf{Protein MSA Conservation Scoring}\\
\large BMI3 Algorithms ICA, Group 1}
\author{}
\date{}
\begin{document}
\maketitle
\vspace{-2.4em}

\section{Input and output}
The software scores positions in rectangular protein multiple sequence alignments (MSAs). Supply aligned FASTA with unique identifiers. The parser uses 20 canonical residues. Dash and dot count as gaps. Ambiguous symbols contribute to nongap coverage and carry no residue mass. Sites with more than 30\% gaps or no canonical mass have blank scores. TSV rows preserve all one-based alignment coordinates. All scores increase with conservation. The segment TSV gives inclusive boundaries, raw-score means, variances and usable-site counts.

\section{Algorithms}
The reference reproduces Johansson and Toh's relative von Neumann entropy (RVNE). We compute canonical column probabilities \(p\) with Henikoff sequence weights. The BLOSUM62 \texttt{qij} pair frequencies supply background marginals \(q\) and a diagonal-normalized similarity \(C\). With \(P=\mathrm{diag}(\sqrt p)\), \(Q=\mathrm{diag}(\sqrt q)\) and \(A(s)=C^{\circ(1-s)/s}\), define
\[
  \rho=P A(s_1)P,\quad \sigma=Q A(s_2)Q,\quad
  D=\operatorname{tr}(\rho\log_{20}\rho)-\operatorname{tr}(\rho\log_{20}\sigma).
\]
Here \(A(0)=I\) and \(A(1)=\mathbf 1\mathbf 1^\top\). The background interpolation weight is \(r\). Defaults are \((r,s_1,s_2)=(0.8,0.8,0.8)\). The reference computes the eigenspace-overlap expression. The cached trace path reuses \(\log_{20}\sigma\) and requires only observed eigenvalues. The paper's per-site score is \(c_iD_i\), with \(c_i\) the unweighted nongap fraction.

SCAW-RVNE fits segments of constant mean to \(D_i\). Dynamic programming minimizes total weighted squared error plus \(\beta V\) per segment, where \(V\) is the full profile's weighted variance. Weighted prefix sums give interval statistics. Centered windows are clipped at fitted boundaries. We select a radius that minimizes \(V_{\rm local}/V+\lambda/C_{\rm local}\), with \(C_{\rm local}\) the local coverage sum. The final score is the selected local mean times \(c_i\). Candidate radii and length constraints are configurable.

For \(N\) sequences, \(L\) sites and \(K=20\) residues, profiles cost \(O(NL+LK)\) and single-configuration scoring costs \(O(LK^3)\). Background caching and batching reduce constant factors. Partitioning costs \(O(LW)\), with \(W=L\) by default or a user-selected maximum segment length. Adaptive scoring costs \(O(LH)\) for \(H\) candidate radii.

\section{Installation and scoring}
The recorded run used Python 3.13.5 and NumPy 2.5.2. Run from \texttt{code/rvne} in the BMI3 repository. These commands use PowerShell. In a POSIX shell use \texttt{export PYTHONPATH=src} in place of the second line.
\begin{verbatim}
python -m pip install -r requirements.txt
$env:PYTHONPATH = "src"
python -m unittest discover -s tests
python -m msa_conservation.cli score --help
\end{verbatim}
\newpage
The following PowerShell command scores the example alignment and exports a segment table.
\begin{verbatim}
python -m msa_conservation.cli score --method scaw-rvne `
  --input data/examples/small_alignment.fasta `
  --output results/sites.tsv --segments-output results/segments.tsv
\end{verbatim}
Use \texttt{rvne} for the paper path, \texttt{rvne-optimized} for the cached path, or \texttt{uniform}, \texttt{gaussian}, \texttt{adaptive} and \texttt{segmentation} for individual contextual comparisons. Pass \texttt{--settings} with the frozen benchmark JSON to use retuned settings, or specify \texttt{--r}, \texttt{--s1} and \texttt{--s2}. The README lists the simulator and other context flags.

For the supplied example, the first three rows have the following values with default SCAW-RVNE settings.
\begin{center}
\begin{tabular}{rrrrr}
\toprule
Position & Raw RVNE & Coverage & Final score & Window size\\
\midrule
1 & 1.5084 & 1.000 & 1.5084 & 1\\
2 & 1.2288 & 1.000 & 1.2288 & 1\\
3 & 1.3219 & 1.000 & 1.2998 & 3\\
\bottomrule
\end{tabular}
\end{center}

\section{Benchmarks}
The simulator sets mutation rates independently of BLOSUM62. Training ranks the 1,331 original RVNE settings. Validation selects global RVNE and context settings and classification thresholds. The held-out evaluation covers eight core scenarios and a 63-family sensitivity sweep. Column comparators include Shannon, KL, Jensen--Shannon, BLOSUM62 severity, ClustalW consensus and the unweighted AL2CO variance index.
\begin{verbatim}
$env:OPENBLAS_NUM_THREADS = "1"
python -m benchmarks.run_rvne_benchmark --mode all --sweep `
  --repeats 4 --output results/rvne
\end{verbatim}
Across the core families, mean Spearman was 0.812 for original RVNE, 0.859 for segmentation alone and 0.846 for retuned SCAW-RVNE. Retuned SCAW-RVNE had mean AUPRC 0.999, compared with 0.953 for original RVNE. These results evaluate the simulated mutation-rate target. The repository documents full results, timing and failure cases. Published catalytic-site AUC and real-family generalisation are outside the measured results.

\section{Sources and contributions}
Johansson and Toh (2010), \href{https://doi.org/10.1142/S021972001000494X}{doi:10.1142/S021972001000494X}. Henikoff and Henikoff (1994), \href{https://doi.org/10.1016/0022-2836(94)90032-9}{doi:10.1016/0022-2836(94)90032-9}. BLOSUM62 joint data provenance and other references appear in the repository.

The supplied project requirements and proposal feedback guided the work. OpenAI Codex assisted implementation, validation, benchmarking and drafting. Named group-member contributions are pending confirmation by the group.
\end{document}
